\iffalse
"Inadecuado funcionamiento de un equipo de cómputo"

Descripción de la problemática: Contar qué situaciones se han estado
presentando, fuentes de consulta con datos que permitan soportar que evidencia
la problemática. Desde cuánto hace que se viene presentando esa situación. 

Pérez (2020), en su estudio presenta que entre un 21\% de equipos de cómputo
tienen fallas por bajo rendimiento en la funcionalidad del sistema operativo,
procesador…

Actores primarios (directos e involucrados: usuario o propietarios del computador quienes lo usan, gestión directa de actividades por medio de estos)
Actores secundarios (jefe. Externas que pueden verse afectadas)

Como cierre…..
Causas y efectos:
\fi
\section{Problemática}

Las situaciones a continuación describen lo que nos referimos cuando se habla
de ``Inadecuado funcionamiento de un equipo de cómputo'' tanto para
dispositivos de sobremesa como móviles:

\subsection{Situaciones relacionadas con el mal funcionamiento de hardware}

\textcite{Herat2020} informó que en 2019 se observó que alrededor del mundo se
desperdiciaron 53.6 millones de toneladas métricas y que se estimó que para el
2023 se tendrán 74.7 millones de toneladas métricas de desperdicios
electrónicos. Y que anualmente se desechan entre 1 y 10 millones de toneladas
métricas en los países industrializados y los que están en vías de desarrollo
desechan alrededor de 500 mil y un millón de toneladas métricas. Esta
acumulación de dispositivos se generan normalmente por una serie de actitudes
que \textcite{Hartl2023} relacionan con la psicología del mercado cuando las
empresas que fabrican diversos productos que con el paso del tiempo resultan
tener más capacidad, más velocidad y más eficiencia: ``Los bienes duraderos
que se estropean al cabo de un tiempo inducen a los consumidores a comprar uno
nuevo, lo que estimula la demanda de este producto.''

Ante esta situación, Kaliyavaradhan et al. (2022) como se cita en
\textcite{Seif2024} encontraron que esta actitud ocurre principalmente por los
acelerados cambios y avances que presenta la tecnología ante varias de nuestras
actividades cotidianas y laborales, cosa que lleva a directamente hacer
reemplazos constantes debido a las incompatibilidades que presentan las
herramientas de software con los componentes y capacidades del hardware: ``A
medida que entran en el mercado nuevos dispositivos con las últimas
prestaciones, los consumidores se deshacen de los más antiguos, lo que aumenta
las tasas de residuos''

Un aparato electrónico que se estropea de forma inesperada, normalmente uno
podría pensar que el problema viene a través del uso, sin embargo, también se
encontró de acuerdo con \textcite{Hartl2023} que también ocurre por
conveniencia empresarial, ya que si alguien nota que su máquina no funciona
como se espera o deja de ser útil, entonces se cambia para mantener a flote la
economía de la empresa que diseña el producto:

\begin{displayquote}
Aunque ingenuamente se podría pensar que la vida útil de un producto depende
normalmente del desgaste de sus materiales más alguna influencia aleatoria,
también puede ser el resultado de una elección empresarial. La obsolescencia
planificada significa que las empresas diseñan sus productos de forma que no
duren mucho.
\end{displayquote}

En esto da sentido a lo que expresan \textcite{Zailani2021} cuando se
encuentran con la problemática de que los usuarios que no son experimentados en
la mantención de hardware y deciden tomar las medidas que solo conocen, pero
que resultan ineficientes ante las situaciones que se presentan cuando un
dispositivo tiende a fallar por algún motivo: 

\begin{displayquote}
A veces, algunos productos informáticos del mercado ofrecen lugares de servicio
oficiales de difícil acceso, por lo que los usuarios se muestran reticentes a
repararlos y, al final, prefieren los lugares de servicio públicos. Sin
embargo, algunos usuarios saben que si reparan su ordenador en un lugar
público, es arriesgado porque no hay información clara sobre los costes de
reparación ni sobre los daños.
\end{displayquote}

\subsection{Situaciones relacionadas con las fallas de ciberseguridad}

Por norma general, algunas veces vemos que ``ciberseguridad'' y ``seguridad
informática'' son sinónimos. Sin embargo, la seguridad informática es el área de
proteger la información de los riesgos que conlleva el entorno en donde se
encuentra y que puede tener otros formatos que no siempre están relacionados
con los equipos de cómputo. En cuanto a la ``ciberseguridad'' tiene una cuestión
similar, solo que se especializa en la protección datos almacenados en sistemas
de cómputo; lo que significa que la ciberseguridad hace parte de la seguridad
informática: ``Se entiende que la ciberseguridad tiene como objetivo
principal la protección de la información digital, que está o es transmitida
entre los todos los dispositivos que se encuentran interconectados, por esto,
la ciberseguridad está en la seguridad informática''
\parencite[e.g.,][]{Mosquera2019} 

\textcite{CanoMartinez2022} comenta que los ciudadanos colombianos
